\paragraph{The integral flags of the mixed seed.}
\label{par:count-two-seed-flags}
Keep the eigenbasis just constructed. To compare the seed with the target we first
need the positions of the seed's ordinary flags in these coordinates.

\begin{lemma}[Residual entries and ordinary flags of the mixed seed]
\label{lem:rank-four-mixed-flags}
Use the eigenbasis of Lemma~\ref{lem:rank-four-seed-diagonalization}.
Let \(r\) be the induced representation in that basis. Write \(\operatorname{supp}_{\rm off}(\bar r)\) for the pairs \((i,j)\), \(i\ne j\), for which \(\bar r_{ij}(g)\ne0\) for some \(g\). Then 
\[
 \operatorname{supp}_{\rm off}(\bar r)\subseteq\{(1,0),(3,2),(3,0),(1,2)\}.
\]
 The \((1,0)\) entry is nonzero
somewhere or both \((3,0)\) and \((1,2)\) are nonzero somewhere.
At every prescribed \(\ell\)-adic place, an integral lower
unipotent symplectic frame for the ordinary flag has \((2,0)\)
entry equal to \(\pm1\) and \((2,1)\) entry in the maximal ideal.
The inertia deletion conditions of the mixed seed hold in this
frame. All entries and comparisons may be read over one finite
coefficient extension.
\end{lemma}

\begin{proof}
\proofstep{Compute the residual entries before and after the other coset.}
For \(g\in G_{Fk_0}\), write \(X(g),Y(g)\) for the
residual upper-right entries of the two rank-two
constituents on the original two cosets. The two
diagonal characters of both constituents are
\(1,\bar\eta\), since \(\eta\) comes from \(\Q\).
Put \(\gamma=b/a\) in the basis change \(P\) above.
Multiplying by \(P^{-1}\) on the left and \(P\) on the
right gives
\[
\begin{array}{c|rrrr}
\text{entry} &(1,0)&(3,2)&(3,0)&(1,2)\\ \hline
\text{value}&
(\bar\gamma X-Y)/2 &-(\bar\gamma X-Y)/2&
(\bar\gamma X+Y)/2 &-(\bar\gamma X+Y)/2 .
\end{array}
\]
All other off-diagonal entries are zero. The diagonal
entries have the marked residual order. On the other
coset in \(G_F\), multiplication by the diagonal matrix
of \(\tau\) preserves the same support and gives the
remaining marked characters. Since \(X\) is nonzero
somewhere, either the \((1,0)\) entry is nonzero at some
element, or both \((3,0)\) and \((1,2)\) are nonzero.
On reducing, the scalar twist therefore has no effect on the tested entries: for its character \(\chi\),
\[
 \overline{\chi(g)r(g)}=\bar\chi(g)\bar r(g)=\bar r(g).
\]

\proofstep{Normalize the eigenlines in each ordering of the local constituents.}
At a place above \(\ell\) the decomposition group is
inside \(G_{Fk_0}\). Its two constituents evaluate
\(r_2\) at the distinct places of \(k_0\) above \(\ell\)
because \(\tau\) exchanges those places. The residual
ordinary lines have opposite characters. Choose an element of the local decomposition group at which
\(\bar\eta\ne1\), and write its matrix as
\(\left(\begin{smallmatrix}a_0&b_0\\c_0&d_0\end{smallmatrix}\right)\).
The lattice of the triangular residual extension has
\(a_0\equiv1\), \(d_0\equiv\eta\), and \(c_0\equiv0\).

For an eigenvector \((p,q)^{\mathsf t}\) of eigenvalue \(\lambda\),
\[
       (\lambda-a_0)p=b_0 q,\qquad
       (\lambda-d_0)q=c_0 p.
\]
If \(\bar\lambda=\bar\eta\), the first multiplier is a unit,
so \(|p|\le|q|\) and the line has a generator \((u,1)\).
If \(\bar\lambda=1\), the second multiplier is a unit and
\(|c_0|<1\), so the generator has the form \((1,w)\) with
\(|w|<1\).

In the first ordering, these lines are \((u,1)\) in
the first constituent and \((1,w)\) in the second.
The largest exponent is \(k+t\) on the first line.
Applying \(P^{-1}\) and normalizing gives first flag
column \((1,\gamma u,-1,\gamma u)^{\mathsf t}\).
The second line is proportional to
\((w,-1,w,1)^{\mathsf t}\).
Write \(v_1=(1,\gamma u,-1,\gamma u)^{\mathsf t}\) for the normalized first column. Subtract its first-coordinate multiple from the second generator and normalize the second coordinate:
\[
 v_2=\frac{(w,-1,w,1)^{\mathsf t}-wv_1}{-1-\gamma uw}.
\]

In the other ordering, the largest exponent is on
the second constituent. The corresponding first flag
column is \((1,-u,1,u)^{\mathsf t}\), and the other
line is proportional to
\((w,\gamma,-w,\gamma)^{\mathsf t}\).
Subtract \(w\) times the first column and divide by
\(\gamma+uw\) for the second flag column.
The ultrametric inequality computes both divisors in the two flag normalizations. Since \(|\gamma|=1\) and \(|uw|<1\),
\[
 |-1-\gamma uw|=1,\qquad |\gamma+uw|=1 .
\]

\proofstep{Complete the pair of columns to a symplectic flag.}
In either case the two columns span a Lagrangian
plane, since the original lines lie in separate
symplectic summands. Write the columns as
\((1,p,q,r)^{\mathsf t},(0,1,s,t)^{\mathsf t}\).
The Lagrangian equation is \(t+ps=q\). Complete them
with \((0,0,1,-p)^{\mathsf t}\) and
\((0,0,0,1)^{\mathsf t}\); the resulting lower
unipotent matrix is symplectic by the same equation.
The first three columns span the orthogonal of the
first line. Thus the frame is adapted to the full
stable flag, with characters of exponents
\(k+t,k,0,-t\) before the scalar twist.

In the two orderings the required entries \(q\) and \(s\) of the normalized matrix satisfy
\[
 q\in\{1,-1\},\qquad |s|<1
\]
in positions \((2,0)\) and \((2,1)\), respectively.
The two constituents' ordinary characters reduce,
in this order, to \(\bar\eta,1,\bar\eta,1\).
They coincide with the marked characters because
\(\overline\Psi\) is trivial on the decomposition
group and the diagonal order was fixed at \(\tau\).

\proofstep{Check the scalar twist on the deleted inertia.}
With \(\vartheta=\epsilon_\ell^{e+1}\omega^{-1}\eta\psi_0\),
the square of the scalar twist is
\[
 \delta^2=\epsilon_\ell^{e+1-k}\omega^{-1}\psi_0\phi_0^{-1}
 \quad\text{on }G_F.
\]
On the deep tower \(\psi_0\) has prime-to-\(\ell\) order.
Above \(\ell\), \(\phi_0\) is one on the decomposition
group and \(\omega\) is the Teichm\"uller part of
\(\epsilon_\ell\). On inertia in the cyclotomic kernel,
\(\psi_0\) reduces to one by the prescribed residual
inertia. Its prime-to-\(\ell\) order makes it one.

The two possible square roots are separated by reduction. Thus
\[
 (\delta-1)(\delta+1)=0,\qquad\delta\equiv1\quad\Longrightarrow\quad \delta=1
\]
at these places. On open inertia where
\(\psi_0=\omega=1\), the exponent of \(\delta^2\) is
\(e+1-k\), which gives the doubled weights used
above. At a prescribed inertia element away from
\(\ell\) with trivial residual labels, the tame
tower condition kills \(\psi_0\), and the prime-to-
\(\ell\) character \(\eta\) is one by reduction. The
quadratic character \(\phi_0\) is unramified on the
tower and \(\epsilon_\ell,\omega\) are unramified away
from \(\ell\). The two constituents are trivial on these inertia elements because the type is \(1\oplus\eta\) and \(\eta=1\) on them. Thus \(\delta^2=1\) there and the same reduction test gives
\[
 \delta^2=1,\quad\bar\delta=1\quad\Longrightarrow\quad\delta=1 .
\]
These are the required inertia conditions for the
integral seed and its reduction.

\proofstep{Choose one finite coefficient extension.}
All these matrices descend to the integers of one
finite extension of \(\Q_\ell\). For the representation,
write \(\overline{\Q}_\ell\) as a countable union of
finite fields over \(\Q_\ell\). Their intersections
with the matrix entries cut out closed subgroups
covering the compact Galois group. By the Baire category theorem at least one of these closed subgroups has nonempty interior. Calling it \(H\), compactness then gives
\[
 H\text{ open in }G,\qquad [G:H]<\infty .
\] Adjoining the entries on finitely
many coset representatives, and the finitely many
flag entries, then gives the desired field.
The displayed unit and strict norm bounds persist
under this coefficient choice. Thus the reduced seed
has the required support and nonzero marked entries
over a common residual coefficient extension.
\label{cor:cohomology-count-two-seed-1}
\end{proof}
